\section{Scope, source assessment, and contributions}\label{sec:source}

\subsection{Exactly which source was analyzed}

The supplied paths were read at ProveIt commit
\begin{center}
\texttt{2a9f8e3ee74a1824d57546f0b1b4f1cef5df93d6}.
\end{center}
This pins the analysis even if the repository later changes. The two inputs
are the vendored readable Lean development \cite{vendor,sernau} and
\emph{Beyond the Riemann hypothesis for Dirichlet $L$-functions:
consequences, transfers and limits} \cite{article}.

The complete Git tree of the vendored \code{grh/} directory is
\begin{center}
\texttt{2e133df63d7f4da6b3059c4195324f8e9187ba73}.
\end{center}
It equals the upstream \code{grh/} tree at
\code{72fead102f663d5c7a407964a13dc0cf9d977207}, as checked against the
upstream GitHub tree response. There are 818 files in that directory and
803 Lean modules below \code{src/}. The local import closure of
\code{Main.lean} contains all 803 modules. The package pins Lean
\code{v4.35.0} and Mathlib
\code{56a005855ea17e0b0958f73ec801ae140b0a1679}.

A lexical audit, after removing nested comments and strings, found no
\code{sorry}, \code{admit}, \code{axiom}, \code{unsafe},
\code{native\_decide}, or \code{implemented\_by} code tokens in this
closure. The deliberately incomplete \code{Challenge.lean} is a separate,
unimported statement-comparison target; its placeholders are not part of
the proof closure. This audit does not by itself validate the proof.

No Lean build or independent kernel replay was run in this session.
Andrew Sutherland's verification page for the same upstream commit reports
successful builds, statement and axiom checks, complete kernel replay,
and independent checker acceptance for the readable and compact editions
\cite{sutherland}. Those are external verification reports. Also, a build
of the ProveIt repository root does not build this separately pinned
vendored package; its own directory and toolchain must be used
\cite{vendor}.

\subsection{The exact mathematical input}

The source statement file imports only Mathlib. Its five exported results
include Mathlib's ordinary Riemann hypothesis and Dirichlet GRH with the
following scope:
\begin{enumerate}
\item Every zero of the Riemann zeta function other than the negative even
integers lies on $\Re s=1/2$, with the conventionally valued pole point
$s=1$ explicitly excluded in the total-function formulation.
\item For every positive modulus and every complex Dirichlet character,
$L(s,\chi)\ne0$ when $\Re s>1/2$, apart from the principal character's
pole at $s=1$.
\item Every Dirichlet $L$-function zero in the open critical strip lies on
$\Re s=1/2$. This includes imprimitive characters; there is also a
primitive-character specialization.
\end{enumerate}
These are the declarations in \code{Statements.lean} and \code{Main.lean}
at the pinned revision \cite{sernau}. The conclusion concerns Dirichlet
$L$-functions over $\Q$. It does not quantify over arbitrary automorphic
$L$-functions, nonabelian Artin $L$-functions, or all number-field zeta
functions, and it imposes no simplicity condition on critical-line zeros.
Every use of GRH in this report fits the stated Dirichlet scope.

\subsection{How the proof mechanism leads to these extensions}

The readable reduction compares upper and lower estimates for the same
corrected covolume. Short integral theta-derived forms give an upper
estimate; near-full-rank spanning supplies enough vectors; a product formula
for twisted lattices, with controlled denominators and torsion
specializations, gives a lower estimate. The common height cancels. In the
source notation the resulting estimate is
\[
 m\cdot\mathrm{totalRemainder}\le f a(\log b)^5.
\]
The reduction files then turn this into an eventual one-sided estimate for
a smoothed prime remainder. An analytic pole-exclusion argument gives
nonvanishing to the right of the critical line, and the functional equation
gives the critical-strip conclusion. These roles can be traced directly in
\code{New/Reduction/MainInequality.lean},
\code{New/Reduction/OneSidedBound.lean}, and \code{New/Main.lean}
\cite{sernau}.

The companion article develops two useful interfaces from this setting.
One identifies products of Dirichlet $L$-functions from their prime data;
the other studies information retained by specialization along highly
contacting formal branches. We extend those interfaces rather than
asserting that an arbitrary nonabelian GRH follows from the input.

The first extension starts from the article's quadratic matching theorem
\code{rec:thm:qmatch}. Its open-question paragraph asks for analogous
matching densities for odd-prime-order characters. The quadratic reduction
uses an integral Walsh transform. At order three, the correct local ring is
the Eisenstein ring at $1-\zeta_3$. Complex conjugation adds a symmetry that
forces the second ternary coding weight in one of the two size regimes.
Section~\ref{sec:cubic} proves the exact extremal cubic result and every
equality case. It does not determine the entire nonextremal density
spectrum or the case of every odd prime exponent.

The second extension starts from the equal-length branch kernel, distance,
and pruning results \code{lub:thm:kernel}, \code{lub:thm:distance}, and
\code{lub:thm:dp}. Question \code{lub:q:adversarial} explicitly asks for
unequal observation lengths, higher uniformizer-power coefficient rings,
and faster error-location methods. Section~\ref{sec:unequal-branches}
resolves the unequal-length field case and its optimal recording costs.
The coefficient-ring and fast-decoding parts remain outside the result.

\subsection{Dependency and novelty map}

\begin{center}
\small
\begin{tabularx}{\linewidth}{@{}L{0.32\linewidth}YL{0.23\linewidth}@{}}
\toprule
Result proved here & Contribution & Analytic dependence\\
\midrule
Exact cubic minimum and all extremizers & Sharp extremal answer for every
$r,k$; new extension of \code{rec:thm:qmatch} & Unconditional; standard
Reed--Muller inputs\\[4pt]
Cubic erasure/corruption thresholds & Sharp consequences of that minimum &
Unconditional\\[4pt]
Finite distinguishing primes and extremal counting & General Fourier bound;
exact extremal indicator norms $4/3,16/9$ & Uses supplied Dirichlet GRH\\[4pt]
Unequal-length kernel and jet recovery & Exact formal-series ambiguity and
recovery criterion & Unconditional\\[4pt]
Branch deletion, distance, and budgets & Exact pruning recursion and sharp
clean/maximal-distance costs & Unconditional\\
\bottomrule
\end{tabularx}
\end{center}

The source article describes itself as a merge of 120 reports, containing
869 results across 689 pages, without a fresh independent review of every
claim \cite{article}. This session inspected the source theorem and the
relevant matching and branch-recovery arguments; it did not certify all
869 results. The present proofs are given in full modulo explicitly cited
standard coding and analytic inputs. The coding weight theorems,
$P$-orderings, weighted greedy selection, and the usual distance criterion
are established tools, not claimed discoveries.
